\section{Vehicles}

\subsection{Basilisk}

Basilisks are artillery pieces that can be relaid, reloaded, and fired in a relatively short time. The design of the Earthshaker Cannon allows it to be elevated high enough to fire its shells in an arc towards concealed targets. They can thus be deployed far from the front line, out of the enemy's reach and sheltered from any return fire. Its lack of armour, coupled with an open platform that exposes the crew, makes it too vulnerable to enemy fire during an assault.

The Earthshaker Cannon has the \textbf{Artillery} ability.

\begin{unitstatstable}
15 cm & 4+ & +0 & 5 & Earthshaker Cannon & 150 cm & 2 & 4+ & $-1$
\end{unitstatstable}

\unitimage{basilisk.png}

\unitdivider

\subsection{Caestus Assault Ram}

The Caestus Assault Ram was originally used for boarding actions during battles in space. The revolutionary anti-gravitic technology with which it is equipped was discovered by the techno-archaeologist Arkhan Land. It allowed the craft to be fitted with a powerful anti-gravitic armour, enabling it to operate both during high-speed orbital operations and in atmosphere. Its chassis is extremely resilient, allowing it to drive straight into its target, ramming it, and then disgorge the assault teams it carries.

The Caestus Assault Ram has the \textbf{All-Round Armour}, \textbf{Assault Ram}, \textbf{Flyer}, and \textbf{Transport (2)} abilities.

\begin{unitstatstable}
30 cm & 2+ & +2 & 5 & Siege Melta & 20 cm & 1 & 3+ & $-3$
\end{unitstatstable}

\unitimage{caestus-assault-ram.png}

\unitdivider

\subsection{Dreadclaw}

The Dreadclaw differs from older Drop Pods in that it is able to take off again under its own power after landing. Dreadclaws were designed to return to orbit after landing in order to move units from one battlefield to another, or even to extract troops from a planet's surface. They also use their enormous thruster in combat, burning enemies as they land or move.

The Dreadclaw has the \textbf{Antigrav}, \textbf{Deep Strike (2)}, \textbf{Orbital Strike}, and \textbf{Transport (2)} abilities. Its Flamers use a \textbf{Template (7.5 cm)} and have the \textbf{Reduces Cover ($-1$)} and \textbf{Bombing (1)} abilities.

\begin{unitstatstable}
20 cm & 4+ & +0 & 5 & Flamers & Bomb & Template & 5+ & 0
\end{unitstatstable}

\unitimage{dreadclaw.png}

\unitdivider

\subsection{Arquitor Bombard}

The Arquitor Bombard is a heavy artillery platform designed to operate at the forefront of a Legiones Astartes advance. Fitted with a reinforced chassis and brutal firepower, it is valued where the fighting is most intense. The wide variety of weapons that can be mounted on its chassis gives it a versatile role on the battlefield. It can therefore be used to smash the most stubborn fortifications, or to annihilate masses of enemy infantry and armour.

The Graviton Mortar has the \textbf{Artillery} and \textbf{Graviton} abilities. The Spicula Missile Launcher has the \textbf{Artillery}, \textbf{Battery}, and \textbf{Template (12 cm)} abilities. The Mortar has the \textbf{Artillery} and \textbf{Reduces Cover ($-1$)} abilities.

\subsubsection*{Arquitor Bombard -- Grav Charge}

\begin{unitstatstable}
15 cm & 3+ & +0 & 5 & Graviton Mortar & 120 cm & 1 & Graviton & $-2$
\end{unitstatstable}

\subsubsection*{Arquitor Bombard -- Spicula Missile Launcher}

\begin{unitstatstable}
15 cm & 3+ & +0 & 5 & Spicula Missile Launcher & 120 cm & Template & 3+ & 0
\end{unitstatstable}

\subsubsection*{Arquitor Bombard -- Mortar}

\begin{unitstatstable}
15 cm & 3+ & +0 & 5 & Mortar & 120 cm & 2 & 3+ & 0
\end{unitstatstable}

\unitimage{arquitor-bombard.png}

\unitdivider

\subsection{Fire Raptor}

The Fire Raptor is a variant of the Storm Eagle whose origins date back to a distant, forgotten age. The Fire Raptor's entire transport capacity is sacrificed to carry its ammunition. Its Avenger Bolt Cannon can drown concentrations of enemy infantry under a hail of fire. Each aircraft also has two side turrets manned by an Astartes acting independently and able to fire at targets other than that of its main gun. The Fire Raptor generally serves as a mobile support platform during aerial assaults and is often used in conjunction with Storm Eagles or Stormravens.

The Fire Raptor has the \textbf{All-Round Armour} and \textbf{Flyer} abilities.

\begin{unitstatstable}
\SetCell[r=2]{c,m} 30 cm &
\SetCell[r=2]{c,m} 3+ &
\SetCell[r=2]{c,m} +2 &
\SetCell[r=2]{c,m} 5 &
Avenger Bolt Cannon &
30 cm &
3 &
4+ &
0
\\
& & & &
Twin Heavy Bolters &
30 cm &
2 &
4+ &
$-1$
\end{unitstatstable}

\unitimage{fire-raptor.png}

\unitdivider

\subsection{Land Raider}

The design of the Land Raider precedes the advent of the Imperium by several millennia. Its rediscovery, at the birth of the Imperium, is commonly attributed to the great techno-archaeologist Arkhan Land. The Land Raider's baptism of fire remains a matter of debate among Imperial archaeologists. Some claim that it was at the Siege of Delebrion that a Land Raider first fired its lascannons, while others point to the armour battles of Calysto Platinum and maintain that the great conflict that ravaged that world at the beginning of the Emperor's Great Crusade was its first theatre of operations. An entire Forge World, Anvilus IX, was devoted exclusively to its production, and it spread rapidly across the galaxy thanks to the Emperor's fleets. The Land Raider's popularity among Imperial forces is immense.

The Land Raider has the \textbf{All-Round Armour} and \textbf{Transport (2)} abilities.

\begin{unitstatstable}
\SetCell[r=2]{c,m} 20 cm &
\SetCell[r=2]{c,m} 2+ &
\SetCell[r=2]{c,m} +3 &
\SetCell[r=2]{c,m} 5 &
Twin Heavy Lascannons &
60 cm &
2 &
4+ &
$-2$
\\
& & & &
Heavy Bolter &
45 cm &
1 &
5+ &
$-1$
\end{unitstatstable}

\unitimage{land-raider.png}

\unitdivider

\subsection{Land Raider Explorator}

The Land Raider Explorator is of a very ancient design, but it was rediscovered during the Great Crusade in the vaults of Mars. It is thought that, originally, the Land Raider was used for the exploration and reconnaissance of hostile worlds. Only later, in a more warlike age, was it converted into a powerful war machine. The Land Raider Explorator contains extraordinarily powerful auspex systems, allowing it to analyse the battlefield with extreme precision. It also has jamming systems that prevent nearby enemy Deep Strikes. This tank is thought to be the origin of all modern Land Raider patterns, intended to explore and conquer new worlds. Its sensor systems make it useful both as a command tank and as the lead vehicle of an armoured vanguard.

The Land Raider Explorator has the \textbf{All-Round Armour}, \textbf{Infiltration}, \textbf{Bulldozer}, \textbf{Jamming}, and \textbf{Transport (2)} abilities.

\begin{unitstatstable}
20 cm & 2+ & +3 & 5 & Twin Heavy Lascannons & 60 cm & 2 & 4+ & $-2$
\end{unitstatstable}

\unitimage{land-raider-explorator.png}

\unitdivider

\subsection{Medusa}

Unlike Imperial artillery pieces that lob their shells over enemy walls, the Medusa sends its own directly into the ramparts themselves, turning them to rubble and opening a breach so that the rest of the force can mount the assault. Long sieges are generally accompanied by the dull reports of Medusa siege cannons, firing day and night from entrenched and well-protected positions. Once they have cut a path through the walls and the area has been secured, the Medusas will advance to support the assault, flattening several buildings with each shell until, street after street, the city is nothing but a field of ruins and the enemy has nowhere left to hide.

The Siege Cannon has the \textbf{Artillery} and \textbf{Damages Buildings (AP $-5$/1)} abilities.

\begin{unitstatstable}
15 cm & 4+ & +0 & 5 & Siege Cannon & 150 cm & 1 & 3+ & $-1$
\end{unitstatstable}

\unitimage{medusa.png}

\unitdivider

\subsection{Landing Modules}

Landing Modules are ultimate weapons of terror and surprise, for they plunge into the heart of the enemy lines. Hardly have they landed when their doors open and their occupants disembark to sow death among enemies plunged into confusion. A Landing Module looks in every respect like a survival capsule. It is launched from a ship in low orbit with colossal force. A ring of retro-thrusters then allows the module to steer towards its programmed landing site. Each Landing Module can transport a Space Marine squad or a Dreadnought in spartan conditions. The principal drawback of a Landing Module assault is that heavy vehicles, such as the Vindicator and the Land Raider, are simply too large to be deployed in this way, which can leave the assault force badly under-armed in a prolonged engagement. Besides the extensive deployment of Dreadnoughts, another solution is to use automated weapon systems mounted inside a Landing Module instead of troops. The Deathstorm pattern, in particular, uses rapid-firing Whirlwind missile launchers or Assault Cannons to clear a wide area on landing, opening the way for the Space Marine assault in their wake.

All Landing Modules have the \textbf{Landing Module} ability. Drop Pods have the \textbf{Transport (2)} ability. The Missile Deathstorm uses a \textbf{Template (12 cm)}. The Assault Cannon Deathstorm has the \textbf{Autonomous Defences} ability, and its Assault Cannons have the \textbf{Turret} ability.

\subsubsection*{Drop Pod}

\begin{unitstatstable}
0 cm & 4+ & +0 & -- & -- & -- & -- & -- & --
\end{unitstatstable}

\subsubsection*{Missile Deathstorm}

\begin{unitstatstable}
0 cm & 4+ & +0 & -- & Missile Launcher & 0 cm & Template & 4+ & 0
\end{unitstatstable}

\subsubsection*{Assault Cannon Deathstorm}

\begin{unitstatstable}
0 cm & 4+ & +0 & -- & Assault Cannons & 45 cm & 4 & 5+ & 0
\end{unitstatstable}

\unitimage{landing-modules.png}

\unitdivider

\subsection{Predator}

The Predator Deimos, or more simply called Predator, is one of the oldest Predator tank variants. It is used by the Space Marine Legions in their conquest of the galaxy and is one of the tanks most used by the Legions. The Deimos Predator is distinguished from other Predator variants by the elegant curvature of its turret and sponsons, which are almost circular. This tank has many advantages, such as its thick armour, versatile armament, and speed, which make it very popular among the Legions.

Its Heavy Autocannon has the \textbf{Turret} ability.

\begin{unitstatstable}
\SetCell[r=2]{c,m} 25 cm &
\SetCell[r=2]{c,m} 3+ &
\SetCell[r=2]{c,m} +0 &
\SetCell[r=2]{c,m} 5 &
Heavy Autocannon &
45 cm &
1 &
4+ &
$-1$
\\
& & & &
Heavy Bolter &
45 cm &
2 &
5+ &
$-1$
\end{unitstatstable}

\unitimage{predator.png}

\unitdivider

\subsection{Predator Annihilator}

The Predator Annihilator is a variant of the famous Predator. It is used by the Space Marine Legions to conquer the galaxy. This variant is equipped with multiple lascannons, allowing it to play an anti-vehicle role. Predator Annihilators bring long-ranged firepower that is valued where enemy resistance is strongest.

Its Twin Heavy Lascannon has the \textbf{Turret} ability.

\begin{unitstatstable}
\SetCell[r=2]{c,m} 25 cm &
\SetCell[r=2]{c,m} 3+ &
\SetCell[r=2]{c,m} +0 &
\SetCell[r=2]{c,m} 5 &
Twin Heavy Lascannon &
60 cm &
1 &
4+ &
$-2$
\\
& & & &
Lascannons &
60 cm &
2 &
5+ &
$-1$
\end{unitstatstable}

\unitimage{predator-annihilator.png}

\unitdivider

\subsection{Predator Executioner}

The Predator Executioner is a variant of the famous Predator. This variant has a Heavy Plasma Cannon mounted in a turret. This plasma cannon is a powerful plasma weapon capable of destroying the heaviest armour with disconcerting ease. The Predator Executioner is used in large numbers by the Space Marine Legions, highly valued for its devastating medium-range firepower.

Its Heavy Plasma Cannon has the \textbf{Turret} ability.

\begin{unitstatstable}
\SetCell[r=2]{c,m} 25 cm &
\SetCell[r=2]{c,m} 3+ &
\SetCell[r=2]{c,m} +0 &
\SetCell[r=2]{c,m} 5 &
Heavy Plasma Cannon &
45 cm &
1 &
4+ &
$-3$
\\
& & & &
Heavy Bolter &
45 cm &
2 &
5+ &
$-1$
\end{unitstatstable}

\unitimage{predator-executioner.png}

\unitdivider

\subsection{Predator Infernus}

The Predator Infernus is a variant of the famous Predator. This variant is equipped with a Multi-melta and Flamers, giving it considerable short-range firepower. It is used both to flush opponents out of fortified positions and to destroy enemy armour at short range.

Its Multi-melta has the \textbf{Turret} ability, and its Flamers have the \textbf{Reduces Cover ($-3$)} ability.

\begin{unitstatstable}
\SetCell[r=2]{c,m} 25 cm &
\SetCell[r=2]{c,m} 3+ &
\SetCell[r=2]{c,m} +0 &
\SetCell[r=2]{c,m} 5 &
Multi-melta &
20 cm &
1 &
3+ &
$-2$
\\
& & & &
Flamers &
20 cm &
2 &
4+ &
0
\end{unitstatstable}

\unitimage{predator-infernus.png}

\unitdivider

\subsection{Rhino}

Probably the most common vehicle in the Imperium's contingents, the armoured Rhino troop transport forms the backbone of every Legion's vehicle park. A Rhino enjoys a perfect balance of armour, transport capacity, and speed, and allows its occupants to deploy rapidly, to send squads onto key positions, or to make surgical strikes against enemy lines. The humble Rhino is one of the pillars of the Space Marine Legions. Across the galaxy, the Emperor's servants bring death by fire and steel to His enemies, carried into the heart of the battlefield's hell by these armoured machines. The Rhino has been used since long before the Great Crusade, its origins stretching back into Humanity's distant past, to the dawn of its expansion into space.

The Rhino has the \textbf{Transport (2)} and \textbf{Attached Transport} abilities.

\begin{unitstatstable}
25 cm & 4+ & +0 & Attached & Storm Bolter & 30 cm & 1 & 4+ & 0
\end{unitstatstable}

\unitimage{rhino.png}

\unitdivider

\subsection{Sabre}

Fast, robust, and heavily armed, the Sabre serves the Space Marines as a tank hunter, attacking key enemy targets and destroying them long before they become a threat. The speed of these vehicles allows them to escape counter-attacks and reform to strike the vulnerable flanks of the enemy army, silencing opposing armour. Equipped with a number of advanced weapon systems, ranging from rapid-firing autocannons to exotic beam weapons, few enemies can withstand the firepower of a complete squadron of these fast hunters.

The Neutron Laser has the \textbf{Damage (+1)} ability.

\subsubsection*{Sabre -- Anvilus Autocannon}

\begin{unitstatstable}
25 cm & 3+ & +0 & 5 & Anvilus Autocannon & 45 cm & 2 & 4+ & $-2$
\end{unitstatstable}

\subsubsection*{Sabre -- Neutron Laser}

\begin{unitstatstable}
25 cm & 3+ & +0 & 5 & Neutron Laser & 75 cm & 1 & 4+ & $-3$
\end{unitstatstable}

\unitimage{sabre.png}

\unitdivider

\subsection{Sicaran}

The Sicaran is based on a unique chassis used by no other vehicle. It was designed during the Great Crusade by Ferrus Manus, Roboute Guilliman, and Tech-Priests who used various STC components to build it, notably those of the Rhino and the Land Raider. The result is a fast vehicle with lethal firepower, and a rate of fire and accuracy without equal. Its main armament consists of two Herakles accelerator autocannons, which allow it to fire shells far faster than a standard autocannon. Thanks to its advanced armament, the Sicaran can track and engage moving targets and find the weaknesses in enemy armour with lethal precision.

Its Herakles Autocannon has the \textbf{Turret} ability.

\begin{unitstatstable}
\SetCell[r=2]{c,m} 25 cm &
\SetCell[r=2]{c,m} 3+ &
\SetCell[r=2]{c,m} +1 &
\SetCell[r=2]{c,m} 5 &
Herakles Autocannon &
45 cm &
2 &
4+ &
$-2$
\\
& & & &
Heavy Bolter &
45 cm &
2 &
5+ &
$-1$
\end{unitstatstable}

\unitimage{sicaran.png}

\unitdivider

\subsection{Sicaran Arcus}

The Sicaran is based on a unique chassis used by no other vehicle. It was designed during the Great Crusade by Ferrus Manus, Roboute Guilliman, and Tech-Priests who used various STC components to build it, notably those of the Rhino and the Land Raider. The result is a fast vehicle with lethal firepower, and a rate of fire and accuracy without equal. The Arcus variant was developed to face large concentrations of enemy infantry. Its main armament is an Arcus launcher that sends volleys of high-explosive missiles, and can even threaten enemies entrenched in fortifications.

Choose one firing mode for the Arcus Launcher when it is used. The Arcus Launcher has the \textbf{Turret} and \textbf{Artillery} abilities. The Arcus Launcher -- Frag has the \textbf{Reduces Cover ($-1$)} ability.

\begin{unitstatstable}
\SetCell[r=3]{c,m} 25 cm &
\SetCell[r=3]{c,m} 3+ &
\SetCell[r=3]{c,m} +1 &
\SetCell[r=3]{c,m} 5 &
Arcus Launcher -- Frag &
90 cm &
2 &
4+ &
0
\\
& & & &
Arcus Launcher -- Krak &
90 cm &
1 &
4+ &
$-2$
\\
& & & &
Heavy Bolter &
45 cm &
2 &
5+ &
$-1$
\end{unitstatstable}

\unitimage{sicaran-arcus.png}

\unitdivider

\subsection{Sicaran Omega}

The Sicaran is based on a unique chassis used by no other vehicle. It was designed during the Great Crusade by Ferrus Manus, Roboute Guilliman, and Tech-Priests who used various STC components to build it, notably those of the Rhino and the Land Raider. The result is a fast vehicle with lethal firepower, and a rate of fire and accuracy without equal. The Omega variant's main armament is a Heavy Plasma Cannon that can vaporise enemy armour and even threaten Titans themselves.

Its Heavy Plasma Cannon has the \textbf{Turret} ability.

\begin{unitstatstable}
\SetCell[r=2]{c,m} 25 cm &
\SetCell[r=2]{c,m} 3+ &
\SetCell[r=2]{c,m} +1 &
\SetCell[r=2]{c,m} 5 &
Heavy Plasma Cannon &
45 cm &
2 &
3+ &
$-3$
\\
& & & &
Heavy Bolter &
45 cm &
2 &
5+ &
$-1$
\end{unitstatstable}

\unitimage{sicaran-omega.png}

\unitdivider

\subsection{Sicaran Punisher}

The Sicaran is based on a unique chassis used by no other vehicle. It was designed during the Great Crusade by Ferrus Manus, Roboute Guilliman, and Tech-Priests who used various STC components to build it, notably those of the Rhino and the Land Raider. The result is a fast vehicle with lethal firepower, and a rate of fire and accuracy without equal. The Punisher variant was developed to scythe down large concentrations of enemy troops at medium range. Its main armament is a Punisher Autocannon, able to fire a true deluge of projectiles, cutting down enemy infantry with ease.

Its Punisher Autocannon has the \textbf{Turret} ability.

\begin{unitstatstable}
\SetCell[r=2]{c,m} 25 cm &
\SetCell[r=2]{c,m} 3+ &
\SetCell[r=2]{c,m} +1 &
\SetCell[r=2]{c,m} 5 &
Punisher Autocannon &
45 cm &
3 &
4+ &
0
\\
& & & &
Heavy Bolter &
45 cm &
2 &
5+ &
$-1$
\end{unitstatstable}

\unitimage{sicaran-punisher.png}

\unitdivider

\subsection{Sicaran Venator}

The Sicaran is based on a unique chassis used by no other vehicle. It was designed during the Great Crusade by Ferrus Manus, Roboute Guilliman, and Tech-Priests who used various STC components to build it, notably those of the Rhino and the Land Raider. The result is a fast vehicle with lethal firepower, and a rate of fire and accuracy without equal. The Venator variant was developed to provide long-ranged heavy fire support. Its main armament is a Neutron Beam Laser, a technological marvel from the finest Forge Worlds. It is capable of destroying the most heavily armoured targets.

Its Neutron Beam Laser has the \textbf{Damage (+1)} ability.

\begin{unitstatstable}
\SetCell[r=2]{c,m} 25 cm &
\SetCell[r=2]{c,m} 3+ &
\SetCell[r=2]{c,m} +1 &
\SetCell[r=2]{c,m} 5 &
Neutron Beam Laser &
90 cm &
1 &
4+ &
$-3$
\\
& & & &
Lascannon &
60 cm &
2 &
5+ &
$-1$
\end{unitstatstable}

\unitimage{sicaran-venator.png}

\unitdivider

\subsection{Storm Eagle}

The Storm Eagle has a role similar to that of Thunderhawks. In many respects it is an ideal compromise between manoeuvrability and firepower. It is less heavily armed than Thunderhawks, but more agile. It is often used to transport Terminator squads while using its formidable array of heavy weapons to support them once they have disembarked.

The Storm Eagle has the \textbf{All-Round Armour}, \textbf{Flyer}, and \textbf{Transport (3)} abilities.

\begin{unitstatstable}
\SetCell[r=2]{c,m} 30 cm &
\SetCell[r=2]{c,m} 3+ &
\SetCell[r=2]{c,m} +2 &
\SetCell[r=2]{c,m} 5 &
Twin Heavy Bolters &
30 cm &
2 &
4+ &
$-1$
\\
& & & &
Vengeance Missile Launcher &
30 cm &
2 &
4+ &
$-2$
\end{unitstatstable}

\unitimage{storm-eagle.png}

\unitdivider

\subsection{Termite}

The Termite is a specialised transport vehicle able to deploy a complete squad of warriors onto the battlefield while bypassing enemy fortifications and sentries. Thanks to its drill, it can bore tunnels through the densest materials at a speed comparable to that of surface transports. Once it has reached its target, it emerges at the surface, or even directly into the lower levels of macro-fortifications, scattering the enemy before disgorging its lethal cargo. The Legions have adopted this mode of transport, which differs from their usual combat doctrine, in order to attack the most important fortifications, thereby bypassing their defences.

The Termite has the \textbf{Deep Strike (2)}, \textbf{Attached Transport}, and \textbf{Transport (2)} abilities.

\begin{unitstatstable}
15 cm & 4+ & +1 & Attached & Autocannon & 45 cm & 1 & 4+ & 0
\end{unitstatstable}

\unitimage{termite.png}

\unitdivider

\subsection{Vindicator}

The Vindicator is a siege tank based on a Rhino chassis and mounting a devastating weapon on its hull: the Demolisher Cannon. The idea of fixing a large-calibre cannon onto a Rhino's hull dates back to the beginning of the Great Crusade. Thunderer Cannons were mounted on modified Rhino chassis to provide heavy fire support in dense urban areas, and thus the Vindicator was born. Its thick armour and lethal firepower have made it popular among the Legions, especially in siege wars and urban environments.

The Demolisher Cannon has the \textbf{Reduces Cover ($-3$)} ability. The Vindicator has the \textbf{Bulldozer} ability.

\begin{unitstatstable}
20 cm & 2+ & +1 & 5 & Demolisher Cannon & 45 cm & 1 & 4+ & $-3$
\end{unitstatstable}

\unitimage{vindicator.png}

\unitdivider

\subsection{Vindicator Destroyer}

The Vindicator is a siege tank based on a Rhino chassis and mounting a devastating weapon on its hull: the Laser Destroyer. The idea of fixing a large-calibre cannon onto a Rhino's hull dates back to the beginning of the Great Crusade. Its thick armour combined with the exceptional firepower of its cannon make it a formidable tank hunter. It can withstand a hail of fire while dealing out lethal shots. These strengths have made it popular among the Legions.

The Vindicator Destroyer has the \textbf{Bulldozer} ability.

\begin{unitstatstable}
20 cm & 2+ & +1 & 5 & Laser Destroyer & 75 cm & 1 & 4+ & $-3$
\end{unitstatstable}

\unitimage{vindicator-destroyer.png}

\unitdivider

\subsection{Whirlwind}

The Space Marines' role as a rapid-attack force means they cannot afford to be slowed by artillery pieces in their ranks. They have nonetheless taken into account the importance of lethal supporting fire, particularly against very numerous or very mobile enemies. It is vital that such opponents be neutralised before they can take advantage of the Space Marines' lack of numbers. It is to fulfil this role that the Whirlwind was conceived. The Whirlwind is a modification of the Rhino chassis, a lightly armoured tank armed with a multiple rocket launcher able to fire precise barrages onto enemy positions.

The Whirlwind may use two different munitions. Choose one firing mode when it is used: Vengeance Missiles or Castellan Missiles. The Vengeance Missile Launcher has the \textbf{Artillery} ability. The Castellan Missile Launcher has the \textbf{Artillery} and \textbf{Reduces Cover ($-1$)} abilities. The Whirlwind has the \textbf{Rapid Intervention} ability.

\begin{unitstatstable}
\SetCell[r=2]{c,m} 25 cm &
\SetCell[r=2]{c,m} 4+ &
\SetCell[r=2]{c,m} +0 &
\SetCell[r=2]{c,m} 5 &
Vengeance Missile Launcher &
100 cm &
1 &
4+ &
$-1$
\\
& & & &
Castellan Missile Launcher &
100 cm &
1 &
4+ &
0
\end{unitstatstable}

\unitimage{whirlwind.png}

\unitdivider

\subsection{Whirlwind Scorpius}

Whirlwinds are mobile artillery pieces that can keep pace with Space Marine intervention forces, which cannot afford to be slowed by artillery in their ranks. Their targeting system and speed allow them to fire while following the assault troops. They are particularly useful against very numerous or very mobile enemies. It is vital that such opponents be neutralised before they can take advantage of the Space Marines' lack of numbers. It is to fulfil this role that the Whirlwind was conceived. The Whirlwind is a modification of the Rhino chassis, a lightly armoured tank armed with a multiple rocket launcher.

The Scorpius Missile Launcher has the \textbf{Artillery} and \textbf{Reduces Cover ($-1$)} abilities. The Whirlwind Scorpius has the \textbf{Rapid Intervention} ability.

\begin{unitstatstable}
25 cm & 4+ & +0 & 5 & Scorpius Missile Launcher & 100 cm & 1 & 3+ & 0
\end{unitstatstable}

\unitimage{whirlwind-scorpius.png}

\unitdivider

\subsection{Xiphon}

The Xiphon Interceptor is an atmospheric combat craft of the Legiones Astartes, of an ancient lineage. Its origins are obscure, but it shared a number of its STC components with craft as diverse as the Amhut Voyager and the Thunderbolt. Able to fly in atmosphere and in space, it combines solid firepower with the speed and agility of an interceptor. Its design is complex, however, and its operational range limited, which demands a great deal of logistics to deploy it. The Xiphon is a high-performance machine that requires constant micro-management from its Astartes pilot in order to keep the craft balanced.

The Xiphon has the \textbf{All-Round Armour}, \textbf{Dodge (5+)}, \textbf{Interceptor}, and \textbf{Flyer} abilities.

\begin{unitstatstable}
-- & 4+ & +5 & 5 & Twin Lascannons & 45 cm & 2 & 4+ & $-1$
\end{unitstatstable}

\unitimage{xiphon.png}
